% GENERATED FILE. Edit canonical lessons, then run npm run generate:books.
% canonical-source-hash: fnv1a64:30927fe6ca23112c
% canonical-lessons: HI-C526-fursat, HI-C526-sangrah, HI-C526-kasida, HI-C526-ikka, HI-C526-satranj

\chapter{Free Time}
\label{ch:hi-free-time}
\hlchaptermodality{526}

\begin{chapteropening}
\textbf{By the end of this chapter:} \emph{I can talk about what I do in my free time.}

Complete the last lesson of chapter 526: I can talk about what I do in my free time.
\end{chapteropening}

\section[\texorpdfstring{\hi{फ़ुरसत} (fursat)}{फ़ुरसत (fursat)}]{\hi{फ़ुरसत} (fursat) — free time; leisure}

\label{lesson:HI-C526-fursat}



\begin{quote}
Before the new one: say the Hindi for to plant (a sapling), then the Hindi for to pull up.
\end{quote}

\subsection*{You'll want to know: \hi{फ़ुरसत}}
\textbf{\hi{फ़ुरसत}} — \emph{fursat} — \textquotedblleft{}free time; leisure\textquotedblright{}.

\begin{quote}
\hi{फ़ुरसत} \hi{में} \hi{मैं} \hi{कहानियाँ} \hi{पढ़ता} \hi{हूँ।}
\end{quote}

\begin{quote}
\emph{Fursat meṁ maiṁ kahāniyā̃ paṛhtā hūṁ.}
\end{quote}

\begin{quote}
In my free time I read stories.
\end{quote}

Feminine.

\begin{cousinweb}
\textbf{\hi{फ़ुरसत}} (\emph{fursat}) is usually traced to Arabic \emph{furṣa}, \textquotedblleft{}an occasion, a chance\textquotedblright{}, which reached Hindi through Persian; in Hindi it is time with nothing that must be done.
\end{cousinweb}

\begin{grammarlens}[title={What you've built}]
One more word for describing what you lived and felt.
\end{grammarlens}

\subsection*{Your turn}
\begin{itemize}
\raggedright
  \item \emph{Say it:} \emph{fursat}
  \item \emph{Say it:} \emph{fursat}, once more
  \item \emph{Say it:} \emph{Fursat meṁ maiṁ kahāniyā̃ paṛhtā hūṁ.} — In my free time I read stories.
  \item \emph{Recall:} say the Hindi for to plant (a sapling), then the Hindi for to pull up, then say \emph{fursat} again
\end{itemize}

\begin{tcolorbox}[breakable,colback=teal!4,colframe=teal!35!black,title={Before you move on}]
What does \hi{फ़ुरसत} mean? (\textquotedblleft{}Free time; leisure\textquotedblright{}.) Say it once more.
\end{tcolorbox}
\section[\texorpdfstring{\hi{संग्रह} (saṅgrah)}{संग्रह (saṅgrah)}]{\hi{संग्रह} (saṅgrah) — a collection}

\label{lesson:HI-C526-sangrah}



\begin{quote}
Before the new one: say the Hindi for to pull up, then the Hindi for free time.
\end{quote}

\subsection*{You'll want to know: \hi{संग्रह}}
\textbf{\hi{संग्रह}} — \emph{saṅgrah} — \textquotedblleft{}a collection\textquotedblright{}.

\begin{quote}
\hi{अरुण} \hi{के} \hi{पास} \hi{पुराने} \hi{सिक्कों} \hi{का} \hi{बड़ा} \hi{संग्रह} \hi{है।}
\end{quote}

\begin{quote}
\emph{Aruṇ ke pās purāne sikkoṁ kā baṛā saṅgrah hai.}
\end{quote}

\begin{quote}
Arun has a big collection of old coins.
\end{quote}

Masculine.

\begin{cousinweb}
\textbf{\hi{संग्रह}} (\emph{saṅgrah}) is a tatsama, taken whole from Sanskrit \emph{saṃgraha}, \textquotedblleft{}holding together, gathering\textquotedblright{} (Monier-Williams).
\end{cousinweb}

\begin{grammarlens}[title={What you've built}]
One more word for describing what you lived and felt.
\end{grammarlens}

\subsection*{Your turn}
\begin{itemize}
\raggedright
  \item \emph{Say it:} \emph{saṅgrah}
  \item \emph{Say it:} \emph{saṅgrah}, once more
  \item \emph{Say it:} \emph{Aruṇ ke pās purāne sikkoṁ kā baṛā saṅgrah hai.} — Arun has a big collection of old coins.
  \item \emph{Recall:} say the Hindi for to pull up, then the Hindi for free time, then say \emph{saṅgrah} again
\end{itemize}

\begin{tcolorbox}[breakable,colback=teal!4,colframe=teal!35!black,title={Before you move on}]
What does \hi{संग्रह} mean? (\textquotedblleft{}A collection\textquotedblright{}.) Say it once more.
\end{tcolorbox}
\section[\texorpdfstring{\hi{कसीदा} (kasīdā)}{कसीदा (kasīdā)}]{\hi{कसीदा} (kasīdā) — embroidery (a pattern sewn in thread)}

\label{lesson:HI-C526-kasida}



\begin{quote}
Before the new one: say the Hindi for free time, then the Hindi for a collection.
\end{quote}

\subsection*{You'll want to know: \hi{कसीदा}}
\textbf{\hi{कसीदा}} — \emph{kasīdā} — \textquotedblleft{}embroidery (a pattern sewn in thread)\textquotedblright{}.

\begin{quote}
\hi{मेरी} \hi{नानी} \hi{दुपट्टे} \hi{पर} \hi{सुंदर} \hi{कसीदा} \hi{करती} \hi{हैं।}
\end{quote}

\begin{quote}
\emph{Merī nānī dupaṭṭe par sundar kasīdā kartī haiṁ.}
\end{quote}

\begin{quote}
My grandmother embroiders beautiful patterns on dupattas.
\end{quote}

Masculine.

\begin{cousinweb}
\textbf{\hi{कसीदा}} (\emph{kasīdā}) is usually traced to Persian \emph{kashīda}, \textquotedblleft{}drawn (out)\textquotedblright{}. The threads are drawn through the cloth.
\end{cousinweb}

\begin{grammarlens}[title={What you've built}]
One more word for describing what you lived and felt.
\end{grammarlens}

\subsection*{Your turn}
\begin{itemize}
\raggedright
  \item \emph{Say it:} \emph{kasīdā}
  \item \emph{Say it:} \emph{kasīdā}, once more
  \item \emph{Say it:} \emph{Merī nānī dupaṭṭe par sundar kasīdā kartī haiṁ.} — My grandmother embroiders beautiful patterns on dupattas.
  \item \emph{Recall:} say the Hindi for free time, then the Hindi for a collection, then say \emph{kasīdā} again
\end{itemize}

\begin{tcolorbox}[breakable,colback=teal!4,colframe=teal!35!black,title={Before you move on}]
What does \hi{कसीदा} mean? (\textquotedblleft{}Embroidery (a pattern sewn in thread)\textquotedblright{}.) Say it once more.
\end{tcolorbox}
\section[\texorpdfstring{\hi{इक्का} (ikkā)}{इक्का (ikkā)}]{\hi{इक्का} (ikkā) — an ace (in cards)}

\label{lesson:HI-C526-ikka}



\begin{quote}
Before the new one: say the Hindi for a collection, then the Hindi for embroidery (a pattern sewn in thread).
\end{quote}

\subsection*{You'll want to know: \hi{इक्का}}
\textbf{\hi{इक्का}} — \emph{ikkā} — \textquotedblleft{}an ace (in cards)\textquotedblright{}.

\begin{quote}
\hi{मीरा} \hi{के} \hi{हाथ} \hi{में} \hi{चार} \hi{इक्के} \hi{थे।}
\end{quote}

\begin{quote}
\emph{Mīrā ke hāth meṁ cār ikke the.}
\end{quote}

\begin{quote}
Mira had four aces in her hand.
\end{quote}

Masculine.

\begin{cousinweb}
\textbf{\hi{इक्का}} (\emph{ikkā}) is usually traced to Persian \emph{yakka}, \textquotedblleft{}single, one alone\textquotedblright{}.
\end{cousinweb}

\begin{grammarlens}[title={What you've built}]
One more word for describing what you lived and felt.
\end{grammarlens}

\subsection*{Your turn}
\begin{itemize}
\raggedright
  \item \emph{Say it:} \emph{ikkā}
  \item \emph{Say it:} \emph{ikkā}, once more
  \item \emph{Say it:} \emph{Mīrā ke hāth meṁ cār ikke the.} — Mira had four aces in her hand.
  \item \emph{Recall:} say the Hindi for a collection, then the Hindi for embroidery (a pattern sewn in thread), then say \emph{ikkā} again
\end{itemize}

\begin{tcolorbox}[breakable,colback=teal!4,colframe=teal!35!black,title={Before you move on}]
What does \hi{इक्का} mean? (\textquotedblleft{}An ace (in cards)\textquotedblright{}.) Say it once more.
\end{tcolorbox}
\section[\texorpdfstring{\hi{शतरंज} (śatrañj)}{शतरंज (śatrañj)}]{\hi{शतरंज} (śatrañj) — chess}

\label{lesson:HI-C526-satranj}



\begin{quote}
Before the new one: say the Hindi for embroidery (a pattern sewn in thread), then the Hindi for an ace (in cards).
\end{quote}

\subsection*{You'll want to know: \hi{शतरंज}}
\textbf{\hi{शतरंज}} — \emph{śatrañj} — \textquotedblleft{}chess\textquotedblright{}.

\begin{quote}
\hi{दादा} \hi{जी} \hi{हर} \hi{शाम} \hi{अपने} \hi{दोस्त} \hi{के} \hi{साथ} \hi{शतरंज} \hi{खेलते} \hi{हैं।}
\end{quote}

\begin{quote}
\emph{Dādā jī har śām apne dost ke sāth śatrañj khelte haiṁ.}
\end{quote}

\begin{quote}
My grandfather plays chess with his friend every evening.
\end{quote}

Feminine.

\begin{cousinweb}
\textbf{\hi{शतरंज}} (\emph{śatrañj}) is usually traced to Persian \emph{shatranj}, which goes back to Sanskrit \emph{caturaṅga}, \textquotedblleft{}having four limbs\textquotedblright{}, the four arms of an army, and already \textquotedblleft{}a kind of chess\textquotedblright{} (CDIAL 4596).
\end{cousinweb}

\begin{grammarlens}[title={What you've built}]
One more word for describing what you lived and felt.
\end{grammarlens}

\subsection*{Your turn}
\begin{itemize}
\raggedright
  \item \emph{Say it:} \emph{śatrañj}
  \item \emph{Say it:} \emph{śatrañj}, once more
  \item \emph{Say it:} \emph{Dādā jī har śām apne dost ke sāth śatrañj khelte haiṁ.} — My grandfather plays chess with his friend every evening.
  \item \emph{Recall:} say the Hindi for embroidery (a pattern sewn in thread), then the Hindi for an ace (in cards), then say \emph{śatrañj} again
\end{itemize}

\begin{tcolorbox}[breakable,colback=teal!4,colframe=teal!35!black,title={Before you move on}]
What does \hi{शतरंज} mean? (\textquotedblleft{}Chess\textquotedblright{}.) Say it once more.
\end{tcolorbox}
